\section{Consequences for the manuscript}

\subsection{The scope of the audit and the standing examples}

The manuscript locations below refer to the supplied primary editions of Chapters~XII and~XVI and the XII appendix \citep{AGGrothendieckXII,AGGrothendieckXIIAppendix,AGGrothendieckXVI}. Their manuscript pagination is distinct from the edition's printed-page numbering. Claims about incorrect formulas on XII manuscript pp.~273--274 were additionally checked against the supplied scanned pages.

\begin{claim}{G.1}{Proved}{The distinction needed to classify the manuscript's uses}
In the modern convention, forward homotopy-pushout invariance is available for every fundamental localizer. A statement asserting reverse cancellation for \emph{all} homotopy-pushout squares is equivalent to the conditions of \claimref{B.5}. The weaker assertion that the intersection and union of two sieves are aspherical only requires forward gluing and retract closure to imply that the pieces are aspherical. It does not imply reverse cancellation for maps of such squares.
\end{claim}
\begin{proof}
The first two assertions are \claimref{B.2} and \claimref{B.5}. For the third, let $K=U\cap V$, $X=U\cup V$, with $K\to e$ and $X\to e$ in $\W$. If $\varnothing\to e\in\W$, product stability implies that every $Y\to e$ is in $\W$, so the conclusion holds. Otherwise $K$ is nonempty; choose $k:e\to K$. Two out of three gives $k\in\W$. The map of spans
\[
(U\leftarrow e\to V)\longrightarrow(U\leftarrow K\to V)
\]
with the identity on $U,V$ induces a specified $\W$-equivalence of Grothendieck constructions. Their second homotopy pushout maps by a $\Winfty$-equivalence to $X$. In the first homotopy pushout, both $U$ and $V$ are retracts in $\Hot_{\W}$: collapse the other summand to its indicated basepoint. Since $X$ is aspherical, these retracts are aspherical by strong saturation. The argument uses no reverse implication. The failure of the corresponding assertion about maps is established in \claimref{G.17} below. This settles the companion's row XIIba-A-053, including the empty-intersection case.
\end{proof}

We use $A$ for the completely specified, checked category of \claimref{B.7}: it has $35$ objects and $159$ arrows, its classifying space is connected and integrally acyclic, and its fundamental group is nontrivial. For a category $T$ and a fresh object $v$, $CT=T^{\triangleright}$ denotes the category obtained by adjoining a terminal object $v$. Its arrows are those of $T$, the identity of $v$, and one arrow $t\to v$ for each $t\in T$, with forced compositions. The category $\Sigma T=C_+T\amalg_T C_-T$ has two such terminal cone vertices and no arrows between them. Its two cone subcategories are full sieves. Their nerve union realizes the unreduced suspension of $BT$. Thus $C_\pm A$ and $\Sigma A$ are weakly contractible, while $A$ is not. Whenever an infinite example is preferable, $A$ can be replaced by the one-object category $BH$ of the nontrivial integrally acyclic group constructed in \claimref{B.10}. This replacement has no nonidentity idempotents. Its cone comma categories satisfy $(BH)/v\cong BH$ literally. The statements below that use this alternative require no finiteness hypothesis.

\begin{claim}{G.2}{Refuted}{The global reverse claims at the start of XII \S5}
XII Th\'eor\`eme~2, manuscript p.~78, and its interpretation of reverse excision as a property of ordinary $\mathrm{HOT}$ at pp.~79--80 are false. The forward clauses of W(7) and W(7 bis) hold for $\Winfty$; the reverse clauses, and the cube clause later called Loc~$(6c)$, fail. The conditional implications between those axioms at pp.~76--78 remain valid, with retract closure for the object-to-cube reduction.
\end{claim}
\begin{proof}
In the sieve-cover square
\[
\begin{tikzcd}
A\ar[r]\ar[d]&C_+A\ar[d]\\
C_-A\ar[r]&\Sigma A,
\end{tikzcd}                                                   \tag{G.2.1}
\]
the lower horizontal map is a weak equivalence, whereas the upper horizontal map is not. This refutes W(7 bis) and hence W(7), since \claimref{A.10} proves the implication W(7)$\Rightarrow$W(7 bis). The finite categorical axioms, sieve conditions, and homology calculations are checked in \texttt{finite\_examples.py}; the nontrivial fundamental group and weak contractibility of the suspension are proved in \claimref{B.7}. The group version uses $BH$ in the same square. The source's corollary ``pieces and union aspherical imply intersection aspherical'' is refuted by the same example. Its forward object corollary is \claimref{B.2}, and its other object corollary is \claimref{G.1}. The assertion about the dual derivator does not repair the claim about $\mathrm{HOT}$: changing a derivator to its opposite changes pushouts to pullbacks, whereas passing from sieves to cosieves by taking opposite categories does not. The relevant companion rows are XIIba-A-049--058 and XIIba-A-062.
\end{proof}

\subsection{XII \S5: sieve pushouts and total asphericity}

\begin{claim}{G.3}{Refuted}{Reverse cancellation remains false along genuine Dwyer maps}
The reverse directions asserted conditionally in XII Proposition~6, manuscript pp.~81--83, Th\'eor\`eme~3, pp.~108--109, and footnote~124, pp.~227--228, do not hold for $\Winfty$, even when the attaching sieve is a genuine Dwyer inclusion and all the source's boundary and closure conditions hold. Their forward directions are valid. The full conditional statements are valid when the appropriate reverse axiom is assumed.
\end{claim}
\begin{proof}
Use the actual $\Phi _2$ mapping cylinder of $A\to e$ from \claimref{B.4}, giving a Dwyer inclusion $i:A\hookrightarrow M$, a weak equivalence $M\to e$, and its pushout
\[
\begin{tikzcd}
A\ar[r,hook,"i"]\ar[d,"p"']&M\ar[d,"q"]\\
e\ar[r,"i'"']&Q=e\amalg_A M.
\end{tikzcd}                                                   \tag{G.3.1}
\]
The Dwyer pushout theorem identifies the specified nerve comparison with the homotopy pushout, so $BQ\simeq\Sigma BA$. Therefore $i',q\in\Winfty$ and $i,p\notin\Winfty$. The generated cosieve of $A$ in $M$ retracts onto $A$ by a right adjoint; the same is true after pushout. Consequently the closure inclusions required in Proposition~6 are categorical homotopy equivalences. Every genuine Dwyer map also satisfies the ``$\W$-dwyerienne'' boundary hypothesis used in Th\'eor\`eme~3. The sizes are $|\Ob M|=141$, $|\Fl M|=1149$, $|\Ob Q|=107$, and $|\Fl Q|=867$; the functor, quotient, sieve, and right-adjoint bijections are checked by the program accompanying \claimref{B.8}.

For the conditional proof, replace the two corners by the specified homotopy-pushout model. The forward implications follow from \claimref{B.2}; the reverse implications follow from \claimref{B.5}. In the more general boundary-weight formulation of Th\'eor\`eme~3, its categorical argument reduces to these comparison squares. No property of arbitrary ordinary pushouts has been substituted for a property of Dwyer pushouts. This records XIIb-a-B002--006, XIIba-C022--024, XIIba-C034, and XIIc-a-C043--045. The implication ``a Dwyer map is a $\W$-cofibration'' itself survives for every modern fundamental localizer, by \claimref{A.11}.
\end{proof}

\begin{claim}{G.4}{Refuted}{The one-point-complement converse}
The one-point-complement equivalence at XII, manuscript p.~85, and again Corollaries~4--5, pp.~149--151, is valid under W(7 bis), but false for $\Winfty$. So are reverse formulations in Proposition~7, pp.~84--85, when specialized to a one-point complement. Their direct implication from aspherical comma categories is valid for every fundamental localizer.
\end{claim}
\begin{proof}
Put $X=\Sigma A$, $U=C_+A$, and let $t$ be the other cone vertex. Then $U$ is a full sieve, $X\setminus U=\{t\}$ has only its identity, and $U/t\cong A$. Both $U$ and $X$ are weakly contractible, so $U\to X\in\Winfty$, while $U/t$ is not weakly contractible. For Corollary~5 take $t_0=t$; its hypothesis on all other complementary objects is vacuous. This also shows why membership of a sieve inclusion in $\Winfty$ cannot be replaced by its universal asphericity under changes of the complementary category.

For completeness, the general conditional one-point argument is a homotopy-pushout calculation. For $H:U^{\op}\to\mathbf{Set}$ with $H(u)=\Hom_X(u,t)$, its category of elements is $U/t$. The nerve of the category formed by adjoining $t$ is the homotopy cone of the projection $U/t\to U$, as in \claimref{B.9}. Reverse cancellation therefore identifies $U\to X\in\W$ with $U/t\to e\in\W$. Without that axiom, the forward direction follows either from the same pushout calculation or directly from Theorem~A, since $U/x$ is terminal for $x\in U$. The data and counterexample also apply with $A=BH$. The affected rows are XIIb-a-B012--015, XIIb-b-B050, B054--055, and XIIapp-C060.
\end{proof}

\begin{claim}{G.5}{Refuted}{The total-asphericity gluing argument has a separate error}
XII Proposition~9(b), manuscript pp.~97--98, its retraction-based repair at pp.~98--100, and its conclusion that the outgoing fork is totally aspherical are false even for $\W=\W_R$ with $R$ any nonzero commutative ring. These errors persist when full W(7) is valid. The additional retractions do not make the relevant left Kan extension preserve monomorphisms or intersections.
\end{claim}
\begin{proof}
Let $C$ have objects $x,z,a,b$, identities, and the six nonidentity arrows
\[
\alpha,\beta:x\rightrightarrows z,\qquad
r:z\to a,\quad s:z\to b,\qquad
r\alpha=r\beta:x\to a,\quad s\alpha=s\beta:x\to b.
\]
There are no further nonidentity arrows or nonforced composites. Let $\Psi\subset C$ be the full outgoing fork on $z,a,b$, and write $f:\Psi\hookrightarrow C$. For $c=z,a,b$, the category $c\backslash\Psi$ is nonempty and connected. For $c=x$, it has four objects: $\alpha,\beta,r\alpha,s\alpha$, and four nonidentity arrows, one from each of $\alpha,\beta$ to each of $r\alpha,s\alpha$. It too is connected. Thus $f$ is final in the ordinary covariant sense, and presheaf left Kan extension gives $f_!e_\Psi=e_C$.

The full subcategory $P=\{x,z\}$ is the parallel-arrow category, with $BP\simeq S^1$. The two full subcategories $P\cup\{a\}$ and $P\cup\{b\}$ are its terminal cones, and cover $C$ by sieves. Hence $BC\simeq\Sigma BP\simeq S^2$. Thus $C\to e\notin\W_R$ for $R\ne0$. This is a final image witnessing that $\Psi$ is not totally $\W_R$-aspherical. Its two arrow subcategories and their intersection $\{z\}$ each have a terminal object, so each is totally aspherical for every fundamental localizer. Each arrow subcategory retracts onto $\{z\}$. Indeed, if a category $T$ has terminal object $t$ and $T\to Y$ is final, the composite $e\xrightarrow{t}T\to Y$ is final, forcing the image of $t$ to be terminal in $Y$.

The precise preservation error can be read off at $x$. The subterminal presheaf represented by $z$ on $\Psi$ maps monomorphically into the terminal presheaf represented by $a$ on the arrow category $z\to a$. After the indicated left Kan extension, the relevant map at $x$ identifies the two-element set $\{\alpha,\beta\}$ with a singleton. It is not injective. Retractions of categories do not change this calculation. The complete composition table, all category axioms, the finality commas, and the normalized chain calculation $H_2(C;\mathbb Z)=\mathbb Z$ are checked in \texttt{ag\_audit\_examples.py}. These calculations settle XIIb-a-B045--049, B051--055 and correct any explanation that attributes their failure solely to the false assertion $\Winfty\models\mathrm W(7)$.
\end{proof}

\begin{claim}{G.6}{Proved}{The suspension construction used to disprove cofibrancy survives}
The counterexample mechanism of XII, manuscript pp.~104--107, survives for $\Winfty$ when its auxiliary category $Y$ is nonempty and simply connected. Its conclusion that the resulting $X'$ is weakly contractible exactly when $Y$ is weakly contractible has an independent proof. The analogous appendix calculation at manuscript pp.~15--16 has the advertised noncofibrancy conclusion, but its displayed fundamental-group calculation is false.
\end{claim}
\begin{proof}
In the construction at pp.~104--107, the nerve of $X'$ is the double cone on $BY$. If $BY$ is simply connected, then $\Sigma BY$ is weakly contractible if and only if $BY$ is weakly contractible. One implication follows from suspension. For the other, reduced integral homology of $BY$ vanishes by the suspension isomorphism. A simply connected acyclic CW complex is weakly contractible: otherwise a first nonzero homotopy group would be its first nonzero reduced homology group by Hurewicz. Nerve realizations are CW complexes, so this argument applies.

In the appendix, its $V'_0$ is exactly the four-object category $C$ of \claimref{G.5}, after renaming $x,z,a,b$ as $0',0,a,b$. Thus $BV'_0\simeq S^2$, not $S^1$, and the assertion $\pi_1(V'_0)\simeq\mathbb Z$ at p.~15 is false. Adjoining the two cone vertices gives $Q$ with $BQ\simeq S^3$, so the pushed-out terminal-cone inclusion is nevertheless not a weak equivalence. For the later square, $\overline V\hookrightarrow X=\Phi_n\times\Phi_m$ is a closed Dwyer inclusion and a weak equivalence: $\overline V$ has a terminal object and both factors $\Phi_n,\Phi_m$ have contractible interval realizations. Thus its pushout $Q\to X'$ is a weak equivalence by the \emph{forward} Dwyer theorem. Two out of three, not reverse excision, now proves $\overline p\in\Winfty\Longleftrightarrow q\in\Winfty$. The closed Dwyer adjoint on its generated sieve $\{0,1,2\}^2$ is coordinatewise $0\mapsto0$, $1\mapsto1$, $2\mapsto1$.

The exact six-object, twenty-arrow double cone is checked in \texttt{ag\_audit\_examples.py}: its normalized chain ranks are $(6,14,16,8)$ and its only positive homology is $H_3=\mathbb Z$. The complete finite $n=m=2$ attachment and its adjoint are checked in \texttt{ag\_appendix\_square.py}. Thus the proof can use homology in place of the invalid one-point reverse inference or the incorrect fundamental-group calculation. This distinguishes XIIba-C008, C012, C016 and XIIapp-C060--065 from false conclusions. The stronger sentence ``nonaspherical whenever $\W\subseteq\W_1$'' is false: $C$ is $\W_1$-aspherical, and $Q$ is even $\W_2$-aspherical.
\end{proof}

\subsection{XII \S6: integral squares, smoothness, and presheaves}

\begin{claim}{G.7}{Proved}{What the integral-square arguments actually require}
XII Proposition~13, manuscript pp.~125--126, its cube corollary at pp.~127--128, the homotopy-cocartesian formulations at pp.~132--135, and their proper-integral duals are valid with their stated full reverse assumptions. Their forward clauses, the construction of the integral models, and the comparison with the closed-cover model require no reverse axiom. Their reverse clauses are false for $\Winfty$.
\end{claim}
\begin{proof}
For a span $X_1\leftarrow X_0\to X_2$, the smooth and proper Grothendieck constructions have specified weak equivalences from the corresponding double mapping-cylinder nerve model. The branch inclusions and the natural transformations of the manuscript exhibit the same coherent cocone. Objectwise $\W$-equivalences induce $\W$-equivalences of these constructions by \claimref{B.2}. This proves the forward cube assertion without reverse excision. The closed-cover replacement is made by the comma categories specified in the text; their branch inclusions have adjoints. It therefore preserves each of the four vertex conditions by actual weak equivalences, not merely by an abstract square isomorphism in $\Hot_\W$.

Under the reverse axiom, opposite-edge cancellation is \claimref{B.5}, and the three cube clauses are \claimref{A.12}. For $\Winfty$, apply the integral construction to the two cones of $A$. The comparison to (G.2.1) is a weak equivalence, so the failed reverse edge and reverse cube conditions transfer by two out of three. This settles XIIbb-A028--035, A043, XIIb-b-B001, B005--006. Footnote~85, p.~126, correctly removes the reverse assumption from the forward edge implications; it does not establish either converse.
\end{proof}

\begin{claim}{G.8}{Refuted}{The reverse smoothness theorem has an explicit finite counterexample}
The reverse direction of XII Th\'eor\`eme~4(a), manuscript pp.~146--147, and its local version in Corollary~6(b), p.~157, are false for $\Winfty$ when the strong W(7 bis) hypothesis is removed. The forward theorem and Corollaries~1--3, pp.~147--149, are valid without reverse excision. The reverse assertions remain valid under their stated full W(7 bis) hypothesis.
\end{claim}
\begin{proof}
For a category $P$, let $T(P)$ be $P$ with a new terminal object $t$, and define $p_P:T(P)\to[1]$ by $P\mapsto0$, $t\mapsto1$. The lifting category for the arrow $0\to1$ at $t$ is precisely $P$. All other lifting categories have terminal objects. Hence $p_P$ is $\W$-smooth if and only if $P$ is $\W$-aspherical.

Take $P_0=A$, $P_1=C_+A$, $P_2=C_-A$, and the inclusions $T(P_0)\to T(P_i)$ that identify the new terminal objects. Form the smooth integral $I$ of
\[
T(P_1)\longleftarrow T(P_0)\longrightarrow T(P_2),
\]
over the base $1\to0\leftarrow2$, and let $\varphi:I\to[1]$ be induced by the three $p_{P_i}$. Explicitly its objects are $(i,x)$, and an arrow $(i,x)\to(0,y)$ for $i=1,2$ is an arrow $x\to\alpha_i(y)$ in $T(P_i)$; arrows within a fibre are its original arrows, with forced composition. There are no arrows between fibres $1$ and $2$. At the terminal object of either outside fibre, the exceptional lifting category for $0\to1$ is $P_1$ or $P_2$, respectively. At the central terminal object, it is the smooth integral of $C_+A\leftarrow A\to C_-A$, whose classifying space is $\Sigma BA$ and hence is weakly contractible. At any object mapping to $0$, and for identity arrows of $[1]$, the lifting category has a terminal object. Thus $\varphi$ is $\Winfty$-smooth, whereas $p_{P_0}$ is not.

The program \texttt{ag\_audit\_examples.py} verifies the entire $110$-object, $1051$-arrow category, all compositions, and the three exceptional lifting categories. The central lifting category has $107$ objects and $867$ arrows and is identified exactly with the indicated integral. Its simple connectivity follows from the double-cone model; the program independently checks integral acyclicity. The same construction works with $BH$ in place of $A$.

For the positive direction, the lifting categories for a smooth integral are integrals of the three lifting categories. When the three are aspherical, their integral is aspherical by forward gluing. For the reverse direction under full W(7 bis), the outside lifting categories are first recovered at the outside objects. Then the equivalence ``the integral and both outside terms are aspherical if and only if the central term is aspherical'' is exactly reverse excision. This gives the required comparison at each arrow of the base. The diamond printed at p.~156 is not strictly commuting in the displayed orientation; its natural transformation supplies the coherent cocone, and the integral comparison is the legitimate homotopy-pushout square. This treats XIIb-b-B035--038, B042--046, XIIbb-C011--020, including the incomplete sentence of footnote~108.
\end{proof}

\begin{claim}{G.9}{Refuted}{The presheaf theorem's converses fail, while its pushout assertion survives}
For a small category $X$, a pushout square of presheaves on $X$ with one leg a monomorphism gives the homotopy-pushout comparison asserted in XII Th\'eor\`eme~5, manuscript p.~168, for every modern fundamental localizer. The two forward equivalence implications hold. The two converses stated there under strong W(7 bis) are false for $\Winfty$ without that hypothesis.
\end{claim}
\begin{proof}
For a presheaf $F:X^{\op}\to\mathbf{Set}$, there is a natural formula
\[
 N(X/F)_n=\coprod_{x_0\to\cdots\to x_n}F(x_n).
\]
All earlier elements in an element-category simplex are obtained by restriction from its last one. Thus the functor $F\mapsto N(X/F)$ preserves colimits degree by degree and sends monomorphisms to simplicial monomorphisms. Applying it to the given presheaf pushout gives a simplicial pushout with a cofibration leg. Its comparison from the homotopy pushout is a weak equivalence, and therefore a $\W$-equivalence. This proves the claimed homotopy statement for an arbitrary second leg, including legs that identify elements. Forward homotopy-pushout invariance gives the two forward equivalence implications. No reflection along pushout is used.

For both converses take $X=\Delta$, $K=NA$, the simplicial cone $CK=K\star\Delta^0$, and the pushout $CK\amalg_K CK$. Both attaching maps are monomorphisms. The category of elements $\Delta/K$ has objects $([n],\sigma)$ with $\sigma\in K_n$ and arrows $([n],\sigma)\to([m],\tau)$ the order-preserving maps $\theta:[n]\to[m]$ with $\sigma=\theta^*\tau$. Its realization is naturally weakly equivalent to $|K|$. Consequently the two functors $\Delta/K\to\Delta/CK$ are not weak equivalences, whereas both $\Delta/CK\to\Delta/(CK\amalg_K CK)$ are. This is a completely specified presheaf counterexample obtained from the checked finite category $A$; the category of all simplices has infinitely many degenerate simplices, which is allowed in the theorem. Under full reverse excision the converses follow from the now-established homotopy-pushout assertion and \claimref{B.5}. This settles XIIbb-C048--052 and confirms the one-way correction in footnote~111.
\end{proof}

\begin{claim}{G.10}{Refuted}{Cofibrancy does not recognize the manuscript's boundary condition}
The assertion at XII, manuscript p.~111, that a $\W$-cofibration sieve with one-point complement is necessarily $\W$-dwyerienne is false. So is its asserted consequence that cofibrancy after every change of the complementary category forces that boundary condition. These failures occur for $\Winfty$, for all $\W_n$, and for $\W_R$ with $R\ne0$, including localizers satisfying full W(7).
\end{claim}
\begin{proof}
Let $U=\{a\}$ and let $X$ have objects $a,t$, their identities, and exactly two arrows $\alpha,\beta:a\to t$. The sole boundary weight is the two-element presheaf on $e$. Its category of elements is the discrete two-point category, which is not aspherical for any localizer listed in the statement. Thus the weight fails even the identity test for total asphericity, and $U\hookrightarrow X$ is not $\W$-dwyerienne.

We prove the stronger cofibrancy assertion with all its quantifiers. Change the complement to an arbitrary small category $T\to e$. The resulting category $X_T$ has objects $a$ and the objects of $T$, the arrows of $T$, and two arrows $a\to t$ for each $t\in T$; postcomposition preserves the labels $\alpha,\beta$. For a functor $e\to C$ choosing $c\in C$, its pushout along $U\hookrightarrow X_T$ is the collage $P(C,c,T)$ with parts $C,T$ and crossing arrows
\[
\Hom_{P(C,c,T)}(y,t)=\Hom_C(y,c)\times\{\alpha,\beta\}.
\]
The crossing-arrow category is exactly $(C/c)\times\{\alpha,\beta\}\times T$. A functor $f:C\to D$ under $e$, with $f(c)=d$, induces $C/c\to D/d$; both categories have terminal objects, so this functor is in $\Winfty$. Products with the fixed category $\{\alpha,\beta\}\times T$ preserve weak equivalences. The gluing comparison proved in \claimref{G.19} therefore expresses $P(f):P(C,c,T)\to P(D,d,T)$ as the induced map of homotopy pushouts of a diagram whose three vertex maps are $f$, this weak equivalence, and $\id_T$. If $f\in\W$, forward gluing gives $P(f)\in\W$. This is precisely the defining condition that $U\hookrightarrow X_T$ be a $\W$-cofibration. When $T=\varnothing$, it is an identity inclusion and the same proof applies.

The four-arrow category and its compositions are included in \texttt{ag\_small\_checks.py}; the universal pushout proof, not a finite sample of $C,D,T$, proves cofibrancy. This verifies the relevant TOTASPH counterexample independently and corrects XIIba-C037--038. The source's narrower question about \emph{trivial} cofibrations is different: the displayed inclusion is not a weak equivalence for $\Winfty$ or for nonzero homology localizers, since its nerve is a circle.
\end{proof}

\subsection{XII \S7 and the previously unaudited end of \S8}

\begin{claim}{G.11}{Proved}{The actual cylinder and cofibration formalism survive}
The factorizations and enough-cofibration results at XII, manuscript pp.~228--234, and Proposition~16 at pp.~231--232 are valid for every modern fundamental localizer. They use the product of the endpoint inclusion with the source category. They do not require the pushout-product stability sought at pp.~114--115. The reverse clauses in footnote~124 are the false extra assertions isolated in \claimref{G.3}.
\end{claim}
\begin{proof}
For $n\geq2$, the construction is
\[
Z(f)=Y\amalg_{X\times\{0\}}(X\times\Phi_n),\qquad
X\xrightarrow{x\mapsto(x,2n)}Z(f)\xrightarrow{p_f}Y.
\]
The map $p_f$ is induced by $f\operatorname{pr}_X$ and $\id_Y$. The endpoint product $X\times\{0,2n\}\hookrightarrow X\times\Phi_n$ has the Dwyer right-adjoint description of \claimref{B.4}; its pushout gives the required Dwyer inclusion of the other endpoint. The inclusion $Y\to Z(f)$ is a pushout of a trivial Dwyer map and is in $\Winfty$; its retraction $p_f$ is therefore in $\Winfty$. This proves the factorization without any cancellation in the reverse direction. Brown's lemma, with the premises verified in \claimref{A.8}--\claimref{A.11}, gives the relative pushout-invariance and couniversal-equivalence assertions of Proposition~16.

By contrast, the map $(U\times\Phi_n)\cup(X\times\{0,2n\})\to X\times\Phi_n$ sought at p.~114 is a pushout-product. The finite counterexample and general obstruction are \claimref{H.6}. Thus no inference in this audit uses that desideratum to justify the actual cylinder. The distinction settles XIIc-a-C046, C050, C053 and C062, as well as the issue raised in the prompt after TOTASPH.
\end{proof}

\begin{claim}{G.12}{Proved}{Pasting homotopy-pushout squares does not use reverse excision}
XII Proposition~17, manuscript pp.~234--235, and Proposition~18, pp.~236--239, hold for every modern fundamental localizer. In particular, if the left square in a rectangle is $\W$-homotopy cocartesian, the right square is $\W$-homotopy cocartesian if and only if the outer rectangle is. This is cancellation of \emph{pushout squares}, not reflection of equivalences along pushout.
\end{claim}
\begin{proof}
For Proposition~17, let $\mathcal X_Q$ be the fibred integral of a commutative square $Q$ and $\mathcal X$ its subintegral on the first three vertices. The inclusion $i_Y:Y\to\mathcal X_Q$ of the fourth fibre is a weak equivalence: the base has its indicated initial object and the fibred replacement has the corresponding adjunction contraction. The natural transformation $i_Y\varphi_Q\Rightarrow i$ written in the source identifies the two maps in $\Hot$, where $i:\mathcal X\to\mathcal X_Q$ is the inclusion. Two out of three gives $i\in\W\Longleftrightarrow\varphi_Q\in\W$.

For Proposition~18, write the rectangle as $A\to B\to C$ above $A'\to B'\to C'$. The hypothesis gives a specified equivalence
\[
B\amalg_A^h A'\longrightarrow B'.
\]
Taking its homotopy pushout with $C$ over $B$ preserves $\W$ by \claimref{B.2}. Homotopy-pushout associativity identifies its source with $C\amalg_A^h A'$. Thus the comparison from this object to $C'$ is in $\W$ exactly when the comparison $C\amalg_B^h B'\to C'$ is, by two out of three. All comparisons use Grothendieck or cylinder models, whose associativity comparison maps are weak equivalences. This proves both directions.

The proof at p.~239 writes a biconditional between two auxiliary maps $\varphi_1$ and $\psi_1$ using W(7 bis). Under the standing hypothesis $\varphi_1\in\W$, only its forward implication is needed. The remaining comparisons have the contractible comma categories displayed there. Thus its stronger invocation is dispensable; there is no counterexample to Proposition~18 itself.
\end{proof}

\begin{claim}{G.13}{Proved}{The intervening fibre theory has no reverse-cofibration dependency}
The theory of $\W$-fibrations and homotopy-cartesian squares in XII, manuscript pp.~240--261, uses the factorization axiom ``enough $\W$-fibrations'' and its stated terminal-object hypothesis, not W(7 bis) in reverse. In $\Cat$, the all-objects-fibrant hypothesis is automatic from product stability. These results therefore apply whenever Loc~$(7F)$ holds, even when Loc~$(6b)$ fails.
\end{claim}
\begin{proof}
For $S$ fixed, the source's replacement category has objects $p:X\to S$ in $\Fib_\W$. A factorization $X\to\overline X\to S$ makes every object of $\Cat/S$ equivalent to one in this subcategory. The dual of \claimref{A.9} proves that a $\W$-map between fibrant objects over $S$ remains in $\W$ after every base change. This is exactly what is required to define the derived pullback and to compare two choices of replacement. Repeating this argument in the arrow category gives the Cartesian-square and pasting assertions. The terminal map $X\to e$ is in $\Fib_\W$, because its pullback operation is product with $X$ and relative Theorem~A gives product stability. Footnote~130 correctly flags that this last premise needs to be supplied in a general abstract category. No use of $\Cof_\W$ or of the reverse direction of a cocartesian-square axiom occurs in these arguments. In particular $\W_0$, which satisfies $(7F)$ but not $(6b)$, is a valid instance; the exact distinction is \claimref{E.2}.
\end{proof}

\begin{claim}{G.14}{Proved}{The two definitions of homotopy-cocartesian square agree}
XII Proposition~20, manuscript pp.~262--264, is valid for every modern fundamental localizer: the Grothendieck-integral criterion for a homotopy-cocartesian square agrees with the criterion using a $\W$-cofibration replacement and ordinary pushout. Its invocation of the strong relative Theorem~A axiom is legitimate under our conventions, and no reverse excision is required.
\end{claim}
\begin{proof}
For a square $Q=(S,X_1,X_2,X_3)$, factor $S\to X_1$ functorially as a Dwyer map $S\to\widetilde X_1$ followed by $p_1:\widetilde X_1\to X_1$ in $\Winfty$. This gives $\widetilde Q\to Q$ with the identity at the other vertices. The induced map of Grothendieck constructions
\[
\mathcal X_{\widetilde Q}\longrightarrow\mathcal X_Q
\]
is in $\W$, by relative Theorem~A or \claimref{B.2}. Replace the other leg similarly if desired. The canonical comparison from the integral to the ordinary pushout along the Dwyer leg is in $\Winfty$. The comparison from that pushout to $X_3$ belongs to $\W$ exactly when the comparison $\mathcal X_Q\to X_3$ does, by two out of three. This is the asserted equivalence.

The displayed arrow labelled $\lambda$ at p.~263 points from $\mathcal X_Q$ to $\mathcal X_{\widetilde Q}$, opposite to the functor induced by the stated $\widetilde Q\to Q$. Reversing that diagram arrow repairs the typing; the argument uses a zigzag in $\W$, so its conclusion is unchanged. Footnote~134 explicitly identifies the stronger relative axiom. XVI footnotes~15--17, manuscript pp.~23--24, return to this comparison and correctly separate it from reverse excision.
\end{proof}

\begin{claim}{G.15}{Proved}{The opening assertion and the gap of XII \S8 need different repairs}
The opening sentence of XII \S8, manuscript p.~265, reverses the implication proved in Proposition~6 at p.~78. The actual earlier result is W(7)$\Rightarrow$W(7 bis). The proposition stated immediately thereafter, deriving the forward and reverse cube clauses from W(7 bis) and the third clause with retract closure, is valid in the modern convention. Footnote~136 at p.~266 identifies a real missing comparison in its proof, which is supplied by \claimref{A.8}--\claimref{A.12}.
\end{claim}
\begin{proof}
The two degenerate cubes of \claimref{A.10} are the proof at p.~78 and establish its direction without additional premises. At p.~265 the proposed proposition has W(7 bis) as its hypothesis, so it is the converse being investigated. In the modern setting, the implication is \claimref{A.12}. The proof's first step replaces each $f_i$ by a Dwyer cofibration followed by a weak equivalence. To preserve the homotopy-pushout conditions of the two faces, one needs a comparison map between their pushout models in $\W$. Termwise weak equivalences alone do not prove this for an arbitrary ordinary pushout. The functorial Dwyer cylinder and the relative Brown gluing lemma give precisely that comparison. This repairs the step without using the reverse conclusion that it is intended to prove. The historical weaker absolute-Theorem-A issue remains exactly as stated in \claimref{A.16}; it is not hidden by this modern repair.
\end{proof}

\begin{claim}{G.16}{Proved}{The cofibres lemma survives, but the reduction uses reverse detection}
The lemma at XII, manuscript p.~268, is true for every modern fundamental localizer: for a map of homotopy-pushout squares, the square of its four pointed homotopy cofibres is again a homotopy-pushout square. The assertion at pp.~269--270 that the \emph{whole reduction} of W(7) to aspherical vertices has used only the forward clause is false. The proof has already used cofibres to reflect equivalences at p.~267.
\end{claim}
\begin{proof}
Let $f:X_\bullet\to Y_\bullet$ be a coherent map of the two squares. At each vertex let $Z_i=Y_i\amalg^h_{X_i}e$ with the specified cone point. The homotopy colimit of the span $Z_1\leftarrow Z_0\to Z_2$ is the iterated homotopy colimit of the corresponding two-dimensional diagram of spans. Interchanging these two homotopy colimits identifies it with
\[
\bigl(Y_1\amalg^h_{Y_0}Y_2\bigr)
\amalg^h_{X_1\amalg^h_{X_0}X_2}e.
\]
The point diagram has homotopy colimit $e$, since its indexing span is contractible. The specified face comparisons to $Y_3$ and $X_3$ are in $\W$, so forward invariance gives a comparison from this displayed object to $Z_3$ in $\W$. The interchanges can be made in the bar construction, or in the Grothendieck construction with its natural comparison weak equivalences. This proves the lemma without reverse detection.

On the other hand, $f_i\in\W\Longleftrightarrow Z_i\to e\in\W$ is the cofibres criterion of \claimref{B.5}. Its reverse direction fails for $\Winfty$. The source invokes precisely that biconditional at p.~267 and again writes $\lambda\in\W$ ``ou encore'' $U\in\mathrm{Asph}_\W$ in its deployment argument at p.~269. The conclusion that the three \emph{object} assertions are equivalent under suitable excision is separate from the claim that they detect the three \emph{map} assertions. Thus we keep the true lemma and remove the unsupported claim about the weaker hypotheses of the entire reduction.
\end{proof}

\begin{claim}{G.17}{Refuted}{The abstract cube clause and its first corollary are false as printed}
Theorem~3\,$^\circ$ at XII, manuscript p.~274, claims $f_0,f_3\in\W\Rightarrow f_1,f_2\in\W$ from enough cofibrations, a terminal object, and closure of aspherical objects under retracts. It is false. Corollary~1 on the same page, identifying this map assertion with its aspherical-object version, is also false without the missing reverse-detection hypothesis. Both failures already occur in $(\Cat,\Winfty)$.
\end{claim}
\begin{proof}
Consider the following strict map between sieve-cover squares:
\[
\begin{array}{c|cccc}
\text{vertex}&0&1&2&3\\ \hline
X_\bullet&A&A&C_-A&C_-A\\
Y_\bullet&A&C_+A&C_-A&\Sigma A\\
f_\bullet&\id_A&A\hookrightarrow C_+A&\id_{C_-A}&C_-A\hookrightarrow\Sigma A.
\end{array}                                                     \tag{G.17.1}
\]
All arrows in both squares and all four vertical maps are literal inclusions, so the cube commutes strictly. The two squares are homotopy pushouts; indeed they are MV squares. The maps $f_0,f_2,f_3$ are weak equivalences, while $f_1$ is not. The category $\Cat$ has a terminal object, enough $\Winfty$-cofibrations by \claimref{B.4}, and retract closure of weakly contractible objects. Thus every hypothesis actually printed in part~3\,$^\circ$ is met. The program \texttt{ag\_small\_checks.py} checks both sieve covers, every category and inclusion, and all six commuting faces. The weak-equivalence decisions use the certified $A$ and suspension proof.

The aspherical-object implication in Corollary~1 is nonetheless true for $\Winfty$, by \claimref{G.1}. It therefore cannot be equivalent to the failed cube assertion. The printed page, checked against the attached edition PDF, contains no conservativity hypothesis in part~3\,$^\circ$; this is not an artefact of the AI transcription. Adding the reverse axiom repairs the map assertion, as \claimref{A.12} proves.
\end{proof}

\begin{claim}{G.18}{Refuted}{The literal Cartesian version of Cons~2 is false}
At XII, manuscript p.~273, the condition named Cons~2 says ``carr\'e cart\'esien''. Read literally, its claim is false even for $\W_Z$. It therefore is not equivalent to the conservativity axiom stated immediately above it. Replacing ``cart\'esien'' by ``cocart\'esien'', retaining the two cofibration hypotheses, gives the valid intended implication.
\end{claim}
\begin{proof}
The strict pullback square
\[
\begin{tikzcd}
\varnothing\ar[r]\ar[d]&\{0\}\ar[d]\\
\{1\}\ar[r]&{[1]}
\end{tikzcd}                                                    \tag{G.18.1}
\]
has three contractible nonempty vertices and empty fourth vertex. Its two initial arrows are $\W$-cofibrations for every modern fundamental localizer, because adjoining a disjoint fixed summand preserves $\W$. For a nonzero ring $R$, $\varnothing$ is not $\W_R$-aspherical, since $H_0(\varnothing;R)=0\ne R$. This verifies every stated hypothesis and refutes the conclusion. All four categories and the empty strict pullback are checked in \texttt{ag\_small\_checks.py}. The word ``cart\'esien'' is also present in the edition PDF, not just in its transcription.

For the corrected pushout clause, a pushout with a cofibration leg is a homotopy pushout. If $Z_1,Z_2,Z_3$ are aspherical, its edge $Z_1\to Z_3$ is in $\W$; reverse cancellation gives $Z_0\to Z_2\in\W$, hence $Z_0$ is aspherical. Cons~1 is exactly detection by the homotopy cone after a cofibration replacement and is already equivalent, in this setting, to the reverse axiom by \claimref{B.5}. Thus no extra false pullback condition is needed.
\end{proof}

\subsection{The appendix's gluing question}

\begin{claim}{G.19}{Proved}{A sieve and its complementary cosieve always give the proposed homotopy pushout}
Let $U$ be a full sieve in a small category $X$, and let $V=X\setminus U$ be the full complementary cosieve. Let $E$ be the full subcategory of the arrow category $\mathrm{Fl}(X)$ on arrows $a:u\to v$ with $u\in U,v\in V$, with source and target functors $s:E\to U$ and $t:E\to V$. The source-to-target natural transformation gives a specified comparison
\[
\hocolim(BU\xleftarrow{Bs}BE\xrightarrow{Bt}BV)\longrightarrow BX,
\]
and it is a weak homotopy equivalence. This proves the conjecture at XII appendix, manuscript pp.~17--18, for every modern fundamental localizer, without reverse excision. The extra reverse tests proposed in the same paragraph are false for $\Winfty$.
\end{claim}
\begin{proof}
We give a category-level proof to specify the coherent comparison. Let $P$ be the fibred Grothendieck construction of the span, using the indexing category with arrows $1\to0\leftarrow2$. Its objects have three types: objects of $U$, objects of $V$, and arrows $a:u\to v$ in $E$. It retains the internal arrows of these three categories. An arrow from $u'\in U$ to $a\in E$ is an arrow $u'\to u$ in $U$, and an arrow from $v'\in V$ to $a\in E$ is an arrow $v'\to v$ in $V$; there are no other inter-fibre arrows. Composition uses $s$ and $t$. Define $\varphi:P\to X$ by $u'\mapsto u'$, $v'\mapsto v'$, and $(u\xrightarrow{a}v)\mapsto v$. It takes the first kind of inter-fibre arrow $b:u'\to u$ to $ab$, the second to its given arrow in $V$, and an internal commutative square in $E$ to its target component. These assignments are a functor precisely because the squares in the arrow category commute.

If $x\in U$, no object mapped into $V$ has an arrow to $x$, so $P/x=U/x$ and has terminal object $\id_x$. Suppose $x\in V$. Define an endofunctor $R$ of $P/x$ by
\[
\begin{aligned}
(u,h:u\to x)&\longmapsto(h:u\to x,\id_x),\\
(a:u\to v,h:v\to x)&\longmapsto(ha:u\to x,\id_x),\\
(v,h:v\to x)&\longmapsto(x,\id_x)\quad\text{in the $V$ fibre}.
\end{aligned}
\]
On an internal $U$ arrow $b$, an internal $E$ square with source component $b$, or an arrow $b$ from the $U$ fibre into the $E$ fibre, it uses the square $(b,\id_x)$. It sends internal $V$ arrows to the identity of its constant value. On an arrow from the $V$ fibre into the $E$ fibre it uses the inter-fibre arrow labelled $\id_x$. The equations defining arrows over $x$ make these assignments composable and functorial. There is a natural transformation $\id_{P/x}\Rightarrow R$: its components are the inter-fibre arrow labelled $\id_u$ on a $U$ object, the square $(\id_u,h)$ on an $E$ object, and $h:v\to x$ on a $V$ object. There is also a natural transformation from the constant functor at the $V$ object $(x,\id_x)$ to $R$, using the inter-fibre arrow $\id_x$ into each $E$ object and the identity on its constant $V$ value. Thus the nerve of $P/x$ is contractible.

Theorem~A gives $\varphi\in\Winfty$. The fibred version of the Grothendieck-construction comparison identifies $BP$ with the displayed homotopy pushout (apply the usual version to opposite categories if desired). Its comparison to $BX$ is exactly $B\varphi$, using the specified source-to-target homotopy. This includes $U=\varnothing$, $V=\varnothing$, or $E=\varnothing$: the same comma-category proof applies and the span becomes the appropriate empty or disjoint union diagram.

Finally, take $X=\Sigma A$, $U=C_+A$, and $V=\{t_-\}$. Then $E\cong A$. Both inclusions $U\to X$ and $V\to X$ are weak equivalences, but neither $E\to U$ nor $E\to V$ is. Thus ``the square is homotopy cocartesian'' cannot be replaced by the reverse equivalences claimed in the question after an edge is assumed weak. These are precisely the false uses of the appendix's footnote~9. The positive comparison proves the stronger $\Winfty$ claim conjectured in footnote~10 and corrects the companion's still-open gluing entry.
\end{proof}

\begin{claim}{G.20}{Proved}{The final conditional theorem is correct after the stated repairs}
XII Theorem~1\,$^\circ$ at pp.~272--273, its equivalence of reverse cube cancellation with conservativity in part~2\,$^\circ$, and Corollary~2 at p.~275 are valid in $\Cat$ for modern fundamental localizers. Corollary~2 must retain its explicit conservativity hypothesis. The same corrected clauses hold in an arbitrary relative category under the precise pushout and factorization premises of \claimref{G.26}, which supplies the compatible finite homotopy-cofibre comparisons from those premises.
\end{claim}
\begin{proof}
The forward cube theorem is \claimref{B.2}. Restricting reverse cube cancellation to the degenerate cubes of \claimref{A.10} gives the single-square conservativity assertion. Its converse and the third cube clause are \claimref{B.5} and \claimref{A.12}, with the necessary retract premise explicit. Hence all three clauses hold together under Corollary~2's hypotheses. In $\Cat$ the functorial cylinder, the all-objects-cofibrant property, and Brown's relative gluing lemma supply the comparisons used in this argument. Footnotes~137--138 ask whether factorization alone supplies the abstract comparison calculus. \claimref{G.26} proves that it supplies the required finite calculus, including the terminal-object comparison; no extra cofibrancy premise on that object is needed. The reference ``XV \S3'' in footnote~137 is to be read as the Brown discussion in XVI \S3. Neither the literal Cartesian Cons~2 nor the unconditional part~3\,$^\circ$ is used in the repaired proof.
\end{proof}

\subsection{XVI \S\S3--7: the complete reverse-dependency register}

\begin{claim}{G.21}{Proved}{The valid and invalid dependencies in XVI \S\S3--4}
The implications in XVI Proposition~3.3, manuscript p.~22, summary~3.5, p.~25, and Corollary~3.7, pp.~26--27, are valid for modern fundamental localizers, with the qualifications on retracts specified in Part~A. They are conditional statements about reversible excision; they do not assert that $\Winfty$ has it. The cofibre calculus in Comment~3.4 and Proposition~3.6, pp.~22--26, and the digression~4.4.15, pp.~66--68, need only forward gluing. The arguments concerning fibration axioms throughout \S4 have no reverse-cofibration premise.
\end{claim}
\begin{proof}
At pp.~16$'$--19, the degenerate cubes give the three formal implications of \claimref{A.10}. At pp.~19--22, the application of Brown's lemma is legitimate once its four premises are checked; \claimref{A.8} supplies the proof the source says it does not know, and \claimref{A.11} supplies the Dwyer premises without a circular use of the desired cube theorem. The conditional reverse implications then give exactly \claimref{A.12}. The unresolved question at p.~23, whether forward single-square excision suffices for the needed cofibre calculus, is answered affirmatively in the modern convention because \claimref{B.2} proves forward cube gluing directly from relative Theorem~A. This does not purport to answer the distinct weaker historical axiom-system question in \claimref{A.16}.

The two questions at pp.~23--24 are the integral criterion and presheaf-monomorphism gluing. They are respectively \claimref{G.14} and \claimref{G.9}. The references there to Th\'eor\`eme~4 and Corollaries~1--3 of XII need only their forward part, as \claimref{G.8} proves. Proposition~3.6 is \claimref{A.11}; its Corollary~3.7 adds exactly the reverse condition $(6C)$. In 4.4.15 the text explicitly removes the unnecessarily inherited W(7 bis) assumption from these forward arguments. Modern strong saturation supplies the premise tentatively mentioned in footnote~35.

The eight fibration formulations and their formal arrows use the dual Brown lemma, path replacements, and at certain historically stated steps sequential-colimit closure or strong saturation; \claimref{A.13}--\claimref{A.17} give their exact premises and modern resolution. None uses reverse excision. The source's $\W_0$ example at pp.~68--70 and footnote~36 is correct: it satisfies $(7F)$ but neither $(6b)$ nor $D_{\Winfty}$. In particular the hoped-for inference from enough fibrations to the fibration axiom is false even with every additional forward, accessibility, and connectedness condition printed there. The exact factorization and classification are in Parts~C and~E.
\end{proof}

\begin{claim}{G.22}{Refuted}{Theorem 5.4's list is false, but its accessibility assertion is true}
XVI Th\'eor\`eme~5.4, manuscript p.~79, is false when its list of axioms includes full Loc~$(6)=(6a)+(6b)$. Its assertions about accessibility of $\Winfty$ and its satisfaction of the other listed axioms remain true. The accessibility proof at pp.~80--83 does not need reverse cofibration excision.
\end{claim}
\begin{proof}
The list includes $(6b)$, refuted by \claimref{G.2}. Its stronger unstated inclusion of $(6c)$ would fail as well by \claimref{G.17}. The remaining listed properties are established in Parts~B, C and~E; in particular $\Winfty$ is the weak-equivalence class of the combinatorial Thomason model category. The class of weak equivalences is accessible, or this may be checked by the countable-extension construction corresponding to the source's convention: begin with countably many objects and arrows in a weak equivalence $A\to B$, close under finite compositions, and successively add witnesses for surjectivity and injectivity of every relative homotopy group. Finite simplicial representatives and finite simplicial homotopies use only finitely many further arrows at each test. There are countably many tests at each countable stage, and countably many stages. Their union is still countable and the induced functor is a weak equivalence. This proves the countable-subcategory extension property stated in Corollary~5.3.2. If there are several components, include witnesses for equality and surjectivity on components first; if both categories are empty, no extension is needed. A weak equivalence with empty target already has empty source.

This argument can be made using the simplicial nerves followed by the accessible fibrant replacement $\mathrm{Ex}^{\infty}$: each finite test factors through a finite stage and a finite part of the categories. It is the same finite-witness mechanism behind the manuscript's path-space and iterated-loop argument. The source's Lemma~5.5 uses the fibration condition to transfer accessibility between maps and aspherical fibres; it does not use $(6b)$. Its unfinished extra hypothesis in part~b is not silently filled in by full excision. The standalone theorem about $\Winfty$ follows by the direct construction above or the model-category accessibility theorem, without that unfinished general lemma. The source's ``$\aleph_0$-accessible'' here means its countable-cardinality convention, not the modern claim that the arrow subcategory is finitely accessible.
\end{proof}

\begin{claim}{G.23}{Proved}{The reverse axiom changes the minimum in XVI \S6}
The fibration theorem XVI~6.1, manuscript pp.~83--85, and Corollary~6.3 have no reverse-cofibration dependency and remain valid; in the modern convention their nontrivial conclusion is $\W=\Winfty$. In contrast, the full stability list proposed in 6.4(b), pp.~87--89, has total localizer as its least solution if it includes $(6b)$. There is no solution satisfying that list and the additionally asserted containment $\widetilde{\W}_\infty\subseteq\Winfty$. If $(6)$ is replaced throughout that list by $(6a)$, the least solution is $\Winfty$.
\end{claim}
\begin{proof}
The proof of 6.1 factors a map as a weak equivalence followed by a reference-perfect map, uses the fibration axiom to make its fibres aspherical, and then iterates based loop spaces. No homotopy pushout or cofibres criterion occurs. Its exact modern version, including the total-localizer case and no accessibility requirement, is \claimref{E.5}. Lemma~6.2's reduction to aspherical objects uses the same factorization and fibre criterion, and so is also independent of the cofibration axiom.

The list in 6.4(b) explicitly includes relative Theorem~A. Therefore modern minimality applies to that list even though the earlier historical definition used absolute Theorem~A. With $(6b)$ present, \claimref{D.4} gives the total localizer; with only $(6a)$ it gives $\Winfty$. Totality contradicts the separately stipulated containment in $\Winfty$. The printed conjecture that this minimum is strictly smaller than $\Winfty$ is therefore false under either proposed reading, for different reasons. Its simple-homotopy suggestion is examined with the correct weaker structures in Part~D. The discussion in 6.4(c), pp.~89--91, that $(7F)$ alone does not identify the canonical path replacement with the needed fibration replacement is correct, and Part~E gives its precise proper-localization meaning. The unresolved question there about $\W_n$ satisfying $(7F)$ is answered affirmatively in Part~C. None of these conclusions about $(7F)$ requires reverse excision.
\end{proof}

\begin{claim}{G.24}{Known}{Accessible localizers give derivators without reverse excision}
For every accessible modern fundamental localizer $\W$, the diagram-localization construction underlying $\mathrm{HOT}_{\W}$ is a strong derivator on all small indexing categories. No assumption of Loc~$(6b)$, Loc~$(7)$, Loc~$(7D)$, or Loc~$(7F)$ is required. In particular the remark at XVI, manuscript pp.~119--120, that $(6b)$ plays no role in obtaining the derivator axioms is correct. This statement uses relative Theorem~A and does not identify it with the weaker historical definition.
\end{claim}
\begin{proof}
By \citet[Theorem~5.2.15, p.~225]{RootCisinski2006}, the category of small categories carries a cofibrantly generated left proper model structure whose weak equivalences are exactly $\W$ and whose cofibrations are the Thomason cofibrations. Its underlying category is locally presentable, so it is combinatorial. Thus \citet[Proposition~1.36, with Theorem~1.34 and Lemmas~1.35--1.37]{AGGroth2013} applies to this very relative category, not just to an abstract equivalent homotopy category. It proves that
\[
J\longmapsto\Ho(\Cat^J)=\Cat^J[(\W^J)^{-1}]
\]
is a strong derivator. With the manuscript's presheaf coefficient convention, put $J=I^{\op}$ and use the corresponding opposite convention on indexing $2$-cells; the displayed values are exactly $\mathrm{HOT}_{\W}(I)$ used in this paper.

Here are the relevant comparisons in that application. Projective and injective model structures on $\Cat^J$ have the same objectwise weak equivalences. For $u:J\to K$, restriction is a right Quillen functor for the projective structures and a left Quillen functor for the injective structures, giving its derived left and right adjoints. The pointwise formulas are computed by homotopy colimits and homotopy limits over the actual comma categories. In the presheaf convention, the left formula at $k\in K$ uses $(k\backslash u)^{\op}$ and the right formula uses $(u/k)^{\op}$. The evaluation and discrete (op)fibration restrictions in these formulas preserve the requisite fibrations or cofibrations: their opposite adjoints are coproducts over the discrete fibres. Consequently the ordinary Kan isomorphism derives to the Beck--Chevalley isomorphism, exactly as in the cited proof. Isomorphisms are detected objectwise because weak equivalences in these diagram model structures are objectwise and saturated; disjoint unions of indexing categories give products of diagram categories. Finally, the model-category cylinder/path rectification for $[1]$ gives the full, essentially surjective partial underlying-diagram functor, which is the strongness assertion of Proposition~1.36.

This verifies the conclusion made in XVI~7.4--7.7, pp.~103--118, under the modern convention and its accessibility hypothesis. The source's tables at pp.~108 and 119 use the forward cofibration theory to construct the required homotopy operations; full Loc~$(6)$ is unnecessarily strong in the first table entry and is corrected to $(6a)$ in the discussion at p.~110. Footnotes~55, 57 and~58 correctly warn that recovery of $\W$ as the localizer of its derivator requires the relative axiom in that argument. We retain it from the outset. We do not assert a model presentation or a derivator existence theorem for arbitrary nonaccessible weak historical localizers. The positive modern theorem here already covers $\Winfty$, all the accessible examples of Part~C, and every accessible reverse-excision example, irrespective of whether it has enough intrinsic fibrations.
\end{proof}

\begin{claim}{G.25}{Proved}{Locations and interpretation of the reverse dependencies}
The following register records the direct and inherited reverse-excision uses in the requested sections. A conditional assertion is recorded as valid only with its actual reverse premise; a redundant use in a proof is distinguished from a false conclusion. ``Forward'' in this table means the precise comparison and premise verified in the cited claim, not unrestricted preservation by ordinary pushout. The supplied XIIc-b excerpt continues through manuscript p.~275, so the register includes the three pages beyond the endpoint mentioned in the prompt.
\end{claim}
\begin{proof}[Verification of the register]
We matched the original transcription page markers to the arguments, checked the companion rows where available, and checked the suspect new printed formulas against the attached XII edition PDF. The executable lexical inventory \path{code/ag_source_inventory.py} locates explicit axiom names and the reverse/\allowbreak conservativity markers; its output is a navigation aid, and the mathematical dispositions below come from the proofs in this paper. The entries of XII \S6 concerning truncation and the intervening purely fibration arguments do not introduce an additional reverse dependency: their separate corrections are in \claimref{E.7} and \claimref{G.13}.
\end{proof}

\small
\begin{longtable}{>{\raggedright\arraybackslash}p{29mm}>{\raggedright\arraybackslash}p{67mm}>{\raggedright\arraybackslash}p{40mm}}
\toprule
Location & Statement or use & Disposition \\
\midrule
\endfirsthead
\toprule Location & Statement or use & Disposition \\\midrule
\endhead
XII 76--80 & Object corollaries, W(7)$\Rightarrow$W(7 bis), Th\'eor\`eme~2 and dual-derivator explanation & \claimref{G.1}--\claimref{G.2}; false reverse for $\Winfty$ \\
XII 81--85 & Proposition~6 cancellation; Proposition~7 and one-point complement & \claimref{G.3}--\claimref{G.4}; forward survives \\
XII 97--100 & Total-asphericity gluing and retraction repair & \claimref{G.5}; false even with full W(7) \\
XII 104--107 & Cone-based obstruction to cofibrancy & \claimref{G.6}; survives by suspension and homology \\
XII 108--114 & Th\'eor\`eme~3, trivial cofibrations, boundary recognition & \claimref{G.3}, \claimref{G.10}; recognition is separately false \\
XII 114--115 & Endpoint pushout-product desideratum & \claimref{G.11}, \claimref{H.6}; cannot replace the actual cylinder \\
XII 125--135 & Proposition~13; integral and homotopy-square reverse comparisons & \claimref{G.7}; conditional reverse valid, false for $\Winfty$ \\
XII 146--157 & Th\'eor\`eme~4 reverse smoothness; Corollaries~4--6; lifting-integral argument & \claimref{G.4}, \claimref{G.8}; finite counterexamples \\
XII 168--169 & Th\'eor\`eme~5 converse implications & \claimref{G.9}; homotopy pushout and forward clauses survive \\
XII 227--234 & Footnote~124 and the cofibre construction & \claimref{G.3}, \claimref{G.11}; remove reverse addenda \\
XII 235--239 & Integral inclusion criterion and pasting Proposition~18 & \claimref{G.12}; reverse use dispensable \\
XII 240--261 & General fibre and Cartesian-square formalism & \claimref{G.13}; no reverse-cofibration use \\
XII 262--264 & Proposition~20 comparison of definitions & \claimref{G.14}; forward only; correct $\lambda$ direction \\
XII 265--267 & Reversed introductory citation, Dwyer replacement, footnote~136, cone detection & \claimref{G.15}; repair comparison; retain reverse detection \\
XII 268--272 & Cofibres square and reduction to aspherical vertices & \claimref{G.16}; lemma true, weaker-reduction claim false \\
XII 273--275 & Cons~1--2, unconditional cube~3\,$^\circ$, Corollaries~1--2 & \claimref{G.17}, \claimref{G.18}, \claimref{G.20}, \claimref{G.26}; two printed errors \\
XII appendix 15--16 & One-point converse and cancellation after closed Dwyer attachment & \claimref{G.6}; conclusions repaired by forward argument \\
XII appendix 17--18 & Open/closed gluing and its proposed reverse tests & \claimref{G.19}; gluing proved, reverse tests false \\
XVI 16$'$--22 & Loc~$(6a,b,c)$, Brown lemma, Proposition~3.3 & \claimref{A.8}--\claimref{A.12}, \claimref{G.21} \\
XVI 23--27 & Integral comparison, presheaf monos, Proposition~3.6, Corollary~3.7 & \claimref{G.9}, \claimref{G.14}, \claimref{G.21} \\
XVI 27--70 & Fibration diagram, 4.4.15 and $\W_0$ example, footnote~36 & \claimref{A.13}--\claimref{A.17}, \claimref{G.21}; reverse not used \\
XVI 71--83 & Accessibility and Th\'eor\`eme~5.4's full list & \claimref{G.22}; list false, accessibility true \\
XVI 83--97 & Fibration theorem, proposed minimum, simple homotopy discussion & \claimref{G.23}, Part~D; full list forces totality \\
XVI 98--123 & Homotopy operations, derivator tables, statement that $(6b)$ is unused & \claimref{G.24}; modern accessible theorem, historical caveat explicit \\
\bottomrule
\end{longtable}
\normalsize

\begin{claim}{G.26}{Proved}{The factorization-only hypotheses supply the required abstract finite calculus}
Let $\mathcal M$ be a category, in a fixed universe large enough for the diagrams under consideration, with all pushouts. Let $\W\subseteq\Fl(\mathcal M)$ contain the isomorphisms and satisfy two out of three. Assume that every arrow factors as an intrinsic $\W$-cofibration followed by a member of $\W$. These are the pushout-dual premises of XII \S7 and the factorization premise at manuscript p.~272. They suffice for a choice-independent notion of homotopy-pushout square, with specified weak comparison maps between its models, invariance under objectwise $\W$-maps, and the compatible finite comparisons needed for the square-of-cofibres lemma at p.~268. If $\mathcal M$ has a terminal object $e$, that lemma holds for the actual cofibres of a map of homotopy-pushout squares whose four vertical maps are cofibrations.

Under these same hypotheses, theorem parts~1\,$^\circ$--2\,$^\circ$ at XII pp.~272--274 are valid after correcting Cons~2 to say \emph{cocartesian}. Corollary~2 is valid with its stated conservativity and object-retract assumptions. No global initial object, cofibrancy or fibrancy of the terminal object, pushout-product axiom, or globally functorial factorization is required. The finite absolute diagrams used here have a least vertex: spans, squares, and their products. We do not assert the existence of arbitrary finite homotopy colimits in $\mathcal M$.
\end{claim}
\begin{proof}
\par\smallskip\noindent\textbf{One-leg gluing.}\quad
Write $C=\Cof_{\W}$. Composition and cobase change preserve $C$, directly by associativity of the corresponding pushout functors, and isomorphisms belong to $C$. By \claimref{A.9}, the maps in $C\cap\W$ are preserved by arbitrary pushout. That claim also proves that a $\W$-map between two $C$-objects under a fixed base stays in $\W$ after any change of that base.

Consider a natural transformation of spans
\[
(A\leftarrow S\to B)\longrightarrow(A'\leftarrow S'\to B')
\]
whose three components $a:A\to A'$, $s:S\to S'$, $b:B\to B'$ are in $\W$. Assume that the corresponding legs $S\to A$ and $S'\to A'$ belong to $C$; no condition is imposed on the other legs. Put $A^*=A\amalg_S S'$. The map $A\to A^*$ belongs to $\W$, since it is the pushout of $s$ along $S\to A$. Therefore the induced map $a^*:A^*\to A'$ lies in $\W$, and both $S'\to A^*$ and $S'\to A'$ belong to $C$. The canonical pushout comparison factors as
\[
A\amalg_S B
\longrightarrow A\amalg_S B'
\cong A^*\amalg_{S'}B'
\longrightarrow A'\amalg_{S'}B'.
                                                        \tag{G.26.1}
\]
The first map is in $\W$ by the definition of $S\to A\in C$, applied to $b$ as a map under $S$. The last is in $\W$ by the relative assertion of \claimref{A.9}, applied to $a^*$ under $S'$. Thus (G.26.1), which is the specified comparison induced by the original transformation, is in $\W$.

\par\smallskip\noindent\textbf{Finite replacements under a fixed object.}\quad
For $T\in\mathcal M$, let $\mathcal M_T=T\backslash\mathcal M$. This category has initial object $\id_T$ and all pushouts, hence all finite colimits. Its full subcategory on the objects $T\to Z$ in $C$ is closed under pushouts along $C$-maps and under the factorizations used below. Finite colimits used in this paragraph are taken in $\mathcal M_T$. The forgetful functor to $\mathcal M$ creates connected colimits: compatibility of an ordinary cocone over a connected diagram makes its induced maps from $T$ coincide. In particular it creates all the nonempty span and product colimits used subsequently. The empty colimit in $\mathcal M_T$ is $T$; no initial object of $\mathcal M$ is inferred.

Let $J$ be a finite poset and $D:J\to\mathcal M_T$ a diagram. A diagram $Q:J\to\mathcal M_T$ is Reedy cofibrant here if, for each $j\in J$, its latching map
\[
L_jQ=\colim_{i<j}Q_i\longrightarrow Q_j
\]
belongs to $C$; an empty latching colimit means $T$. Choose a linear ordering extending the partial order. Having constructed $Q_i$ and $q_i:Q_i\to D_i$ for $i<j$, the maps $q_i$ and the structure maps of $D$ induce a specified map $L_jQ\to D_j$. Factor it as
\[
L_jQ\xrightarrow{c_j}Q_j\xrightarrow{q_j}D_j,
\qquad c_j\in C,\quad q_j\in\W.                         \tag{G.26.2}
\]
The factorization defines all incoming maps to $Q_j$ and their composites. Thus $q:Q\to D$ is a natural transformation, objectwise in $\W$. At a least vertex with value $T$ we may take both arrows to be identities. There are only finitely many choices.

Two elementary colimit facts give the needed control. If $K\subseteq J$ is downward closed and $Q$ is Reedy cofibrant, then
\[
\colim_KQ\longrightarrow\colim_JQ
\]
belongs to $C$. Indeed, add the remaining vertices in a linear ordering; adding $j$ is the pushout of $L_jQ\to Q_j$. In particular all $T\to Q_j$ and $T\to\colim_JQ$ are in $C$. Secondly, an objectwise $\W$-map $Q\to R$ of Reedy-cofibrant diagrams induces a $\W$-map on colimits. Induct on the number of vertices. When a last vertex $j$ is added, both its latching diagram and the previously added downward-closed diagram have fewer vertices, so their comparison maps belong to $\W$ by induction. Apply (G.26.1) to the spans
\[
Q_j\longleftarrow L_jQ\longrightarrow\colim_{i\ne j}Q_i,
\qquad
R_j\longleftarrow L_jR\longrightarrow\colim_{i\ne j}R_i.
\]
Their left legs are the latching cofibrations. The empty case is $\id_T$. This proves colimit invariance without a model-category or lifting assumption.

\par\smallskip\noindent\textbf{Common refinements and change of the initial vertex.}\quad
The independence assertion uses actual comparison maps. Suppose $Q\to D$ and $Q'\to D$ are Reedy-cofibrant replacements in $\mathcal M_T$. Their pointwise coproduct $P=Q\amalg_TQ'$ is Reedy cofibrant, because latching colimits commute with this coproduct and a coproduct of two cofibrations is a composite of two pushouts of cofibrations. There is a natural map $P\to D$. Construct a relative replacement $P\to R\xrightarrow{r}D$ inductively by factoring
\[
P_j\amalg_{L_jP}L_jR
\longrightarrow D_j
\quad\text{as}\quad
P_j\amalg_{L_jP}L_jR\xrightarrow{d_j}R_j\xrightarrow{r_j}D_j,
\qquad d_j\in C,\quad r_j\in\W.                         \tag{G.26.3}
\]
The map $L_jR\to R_j$ is a composite of the pushout of $L_jP\to P_j$ and $d_j$, hence is in $C$. Thus $R$ is Reedy cofibrant. The natural maps $Q\to R$ and $Q'\to R$ are objectwise in $\W$, since their composites with $r$ are the given objectwise $\W$-maps to $D$. Their colimit maps are therefore in $\W$ by the preceding paragraph. These are commuting comparison diagrams over $D$, not an asserted abstract isomorphism after localization.

For a span $D=(D_1\leftarrow D_0\to D_2)$, a cofibrant span model means a span $Q$ with both legs in $C$ and an objectwise $\W$-map $q:Q\to D$. Its root need not equal $D_0$. Rectify that root by putting
\[
Q^*_0=D_0,\qquad Q^*_i=Q_i\amalg_{Q_0}D_0\quad(i=1,2).
                                                               \tag{G.26.4}
\]
The canonical maps $Q\to Q^*$ are objectwise in $\W$, since the two branch maps are pushouts of $q_0\in\W$ along cofibrations. The induced $Q_i^*\to D_i$ lie in $\W$ by two out of three, and the two legs of $Q^*$ are in $C$. Furthermore,
\[
\colim Q\longrightarrow\colim Q^*
\cong(\colim Q)\amalg_{Q_0}D_0
\]
is in $\W$, because $Q_0\to\colim Q$ is a composite of cofibrations. Now $Q^*$ is a Reedy-cofibrant replacement under the fixed root $D_0$, so (G.26.3) compares it to any other such replacement. This proves that every cofibrant span model gives the same test for the membership in $\W$ of its specified comparison to a fourth vertex.

Consequently the definition of a homotopy-pushout square by cofibration replacement is independent of all choices and of an auxiliary object under which the calculation was performed. If one leg of the original span is already in $C$, (G.26.1), applied to a cofibrant span model and that original span, shows directly that the canonical comparison from the model to the ordinary pushout lies in $\W$. Thus an ordinary pushout with one cofibration leg is a homotopy pushout, even when the other leg is a map to a noncofibrant terminal object.

\par\smallskip\noindent\textbf{Product diagrams and the finite interchange.}\quad
For finite posets $J,K$, let $Q:J\times K\to\mathcal M_T$ be Reedy cofibrant. Each fixed row and column is Reedy cofibrant. We check the point that a fixed row is not itself downward closed. For fixed $k\in K$, its latching map at $j$ identifies with
\[
\colim_{J_{<j}\times K_{\leq k}}Q
\longrightarrow
\colim_{J_{\leq j}\times K_{\leq k}}Q
\cong Q(j,k).                                                \tag{G.26.5}
\]
The source is canonically $\colim_{i<j}Q(i,k)$, because $k$ is terminal in $K_{\leq k}$. Both indexing subposets in (G.26.5) are downward closed in the product, and one is included in the other. Thus (G.26.5) belongs to $C$ by the attachment argument. This also treats empty latching subposets. The column argument is identical with the factors interchanged.

The partial-colimit diagram $k\mapsto\colim_{j\in J}Q(j,k)$ is Reedy cofibrant as well: its latching map is
\[
\colim_{J\times K_{<k}}Q
\longrightarrow
\colim_{J\times K_{\leq k}}Q
\cong\colim_{j\in J}Q(j,k),                                  \tag{G.26.6}
\]
again the map from a downward-closed subposet. Ordinary colimits commute, giving the specified canonical isomorphisms
\[
\colim_K\colim_J Q
\cong\colim_{J\times K}Q
\cong\colim_J\colim_K Q.                                     \tag{G.26.7}
\]
Equations (G.26.2)--(G.26.6) show that these are comparisons of valid compatible homotopy-colimit models whenever the relevant diagrams are formed under one fixed object. Different choices are compared by (G.26.3). No pushout-product of maps has entered the construction.

In particular, for a map of spans or squares choose a replacement of its whole diagram, indexed by the original poset times $[1]$, under the value of its least vertex. The two faces are cofibrant span models, including when their initial objects differ from those of the original faces; (G.26.4) justifies that change. If the given transformation is objectwise in $\W$, the replacement's corresponding vertical maps are in $\W$ by two out of three. Equations (G.26.1) and (G.26.7) then give the specified weak comparison of pushout models. The commutative square with the two fourth-vertex comparisons proves invariance of homotopy-pushout squares under objectwise $\W$-maps. More generally, for a map of two homotopy-pushout squares, $f_0,f_1,f_2\in\W$ implies $f_3\in\W$. This is the abstract theorem's part~1\,$^\circ$.

\par\smallskip\noindent\textbf{The actual square of cofibres.}\quad
Assume now that $e$ is terminal. Let $I=[1]\times[1]$, with vertices $0,1,2,3$ in the manuscript's order, and let $J$ be the span poset with arrows $0\to1$, $0\to2$. Given a strict map $f:X\to Y$ of homotopy-pushout squares with each $f_i\in C$, form the actual diagram
\[
D:I\times J\longrightarrow\mathcal M,\qquad
D(i,0)=X_i,\quad D(i,1)=Y_i,\quad D(i,2)=e.                   \tag{G.26.8}
\]
The two $J$-arrows are $f_i$ and the unique $X_i\to e$; the $I$-arrows in the last column are identities. Its least vertex has value $T=X_0$. Choose a Reedy-cofibrant replacement $q:Q\to D$ in $\mathcal M_T$.

Let $I_0=I\setminus\{3\}$, the initial three-vertex span. For $j=0,1,2$ there are natural comparison maps
\[
\lambda_j:\colim_{i\in I_0}Q(i,j)\longrightarrow Q(3,j).
                                                               \tag{G.26.9}
\]
For $j=0,1$, their composites with $q_{3,j}$ are the comparison maps of cofibrant span models to the fourth vertices of $X$ and $Y$. They are in $\W$ by hypothesis and the independence argument. Since $q_{3,j}\in\W$, two out of three gives $\lambda_j\in\W$.

For $j=2$, compare the span $Q(-,2)|_{I_0}$ to the constant span $e\leftarrow e\to e$. Its three comparison maps lie in $\W$, its own legs are cofibrations, and the target legs are identities, hence cofibrations. Equation (G.26.1) proves
\[
\colim_{i\in I_0}Q(i,2)\longrightarrow e\in\W.
                                                               \tag{G.26.10}
\]
The composite of $\lambda_2$ with $q_{3,2}:Q(3,2)\to e$ is this map, so $\lambda_2\in\W$. This is the full justification that the terminal constant span has homotopy pushout $e$; it uses no cofibrancy of $T\to e$.

By (G.26.5)--(G.26.6), the two $J$-diagrams in (G.26.9) are Reedy cofibrant. Colimit invariance therefore sends their objectwise $\W$-map $\lambda$ to a map in $\W$. Put
\[
Z'_i=\colim_{j\in J}Q(i,j).
\]
Using (G.26.7), the resulting map is exactly
\[
Z'_1\amalg_{Z'_0}Z'_2\longrightarrow Z'_3\in\W.               \tag{G.26.11}
\]
The square $Z'$ is Reedy cofibrant by (G.26.6), so (G.26.11) says it is a homotopy-pushout square.

These $Z'_i$ are replacement-row cofibres. The cofibres in the manuscript's lemma are the actual pushouts
\[
Z_i=Y_i\amalg_{X_i}e,
\]
with their specified maps $e\to Z_i$. The natural transformation $q$ induces a natural comparison $Z'\to Z$. At each vertex its map is in $\W$ by (G.26.1), applied to the corresponding row of $Q$ and (G.26.8): the required target leg is precisely $f_i\in C$. Thus the actual square $Z$ is a homotopy pushout by the invariance already established. This supplies the exact square-of-cofibres lemma with actual comparison maps and its actual point data.

\par\smallskip\noindent\textbf{Conservativity and the final abstract theorem.}\quad
For arbitrary $f_i$, first apply (G.26.2) to the whole $I\times[1]$ diagram of $f$, under $X_0$. Its vertical maps $g_i:\widetilde X_i\to\widetilde Y_i$ are in $C$ by (G.26.5); its vertex comparisons to $X,Y$ are in $\W$. Hence $g_i\in\W$ if and only if $f_i\in\W$, and its two faces are still homotopy pushouts. Its actual cone square
\[
Z_i=\widetilde Y_i\amalg_{\widetilde X_i}e
\]
is therefore a homotopy pushout by the preceding paragraph.

Suppose that homotopy-pushout squares are conservative in the source's sense. For a cofibration $g:A\to B$, its ordinary cone square is a homotopy pushout. Thus
\[
g\in\W
\quad\Longleftrightarrow\quad
e\to B\amalg_Ae\in\W
\quad\Longleftrightarrow\quad
(B\amalg_Ae)\to e\in\W.                                    \tag{G.26.12}
\]
The first forward implication also follows directly from \claimref{A.9}; its reverse is conservativity. The second equivalence follows from two out of three, since the composite of the two displayed point maps is $\id_e$.

If $f_1,f_2,f_3\in\W$, then $Z_1,Z_2,Z_3$ are aspherical by (G.26.12). The edge $Z_1\to Z_3$ is in $\W$, so conservativity gives $Z_0\to Z_2\in\W$. Hence $Z_0$ is aspherical and $g_0$, and therefore $f_0$, is in $\W$. Conversely this reverse cube clause implies conservativity by the degenerate cube in \claimref{A.10}; its degenerate face is an ordinary pushout along an identity, hence is a homotopy pushout by (G.26.1). This proves part~2\,$^\circ$.

Conservativity implies Cons~1 and the corrected cocartesian Cons~2 directly. Conversely assume those two conditions. In the preceding cone-square argument, factor the two initial legs of $Z$ as cofibrations followed by $\W$-maps, keeping $Z_0$ fixed. The resulting ordinary pushout $P$ maps to $Z_3$ by a member of $\W$. If $Z_1,Z_2,Z_3$ are aspherical, both replaced branches and $P$ are aspherical. Corrected Cons~2 implies that $Z_0$ is aspherical, and Cons~1 applied to $g_0$ implies $g_0\in\W$. This gives the reverse cube clause and hence conservativity. The literal Cartesian condition remains false, as \claimref{G.18} proves.

Finally assume the source's literal premise that every retract \emph{in $\mathcal M$} of a $\W$-aspherical object is aspherical. In a homotopy-pushout square $Z$ with $Z_0,Z_3$ aspherical, factor its two initial legs as
\[
Z_0\longrightarrow\widehat Z_i\longrightarrow Z_i
\quad(i=1,2),
\]
with first arrows in $C$ and second arrows in $\W$. Let $P=\widehat Z_1\amalg_{Z_0}\widehat Z_2$. Its comparison to $Z_3$ lies in $\W$, so $P$ is aspherical. Push the two legs and this pushout along $Z_0\to e\in\W$. All resulting maps
\[
\widehat Z_i\longrightarrow Z_i^*=\widehat Z_i\amalg_{Z_0}e,
\qquad
P\longrightarrow P^*=Z_1^*\amalg_e Z_2^*
\]
belong to $\W$, since the respective structural maps out of $Z_0$ are cofibrations. Thus $P^*$ is aspherical. Each $Z_i^*$ is an actual retract of $P^*$: use its identity and the composite $Z_{3-i}^*\to e\to Z_i^*$ on the other branch. These maps agree on $e$. Retract closure makes both $Z_i^*$, and then both $Z_i$, aspherical. Applying this object result to the actual cone square and using (G.26.12) proves the third cube clause under conservativity. Together with the forward and reverse clauses this is Corollary~2. It does not prove the unconditional part~3\,$^\circ$ or its Corollary~1, whose counterexample remains \claimref{G.17}.
\end{proof}
